\documentclass[sigplan,anonymous,review,nonacm]{acmart}

\usepackage{amsmath}
\usepackage{booktabs}
\usepackage{algorithm}
\usepackage{algpseudocode}
\usepackage{tabularx}
\algnewcommand{\algorithmicwhen}{\textbf{when}}
\algnewcommand{\When}[1]{\State \algorithmicwhen\ #1 \textbf{do}}
\algnewcommand{\EndWhen}{}

\title{End-to-End Max-Min Fairness for Multimodal LLM Serving}

\begin{document}

\begin{abstract}
Multimodal LLM requests consume substantial service before they reach the
inference engine. A video request may fetch media, parse its container, seek,
decode and resize frames, and construct a model-ready prompt before consuming
vision-encoder, prefill, and generation service. Existing LLM fairness
mechanisms begin inside the inference engine and therefore cannot account for
contention in this upstream preparation stage. Consequently, tenants that
receive similar inference service can experience unequal end-to-end service,
and a tenant with expensive videos can monopolize preparation capacity.

We present an external, work-conserving scheduler that provides unweighted
max-min fairness across tenants at both video preparation and inference
admission. The scheduler maintains a separate virtual-service counter for
each tenant at each stage, dispatches from the least-served backlogged tenant,
and reconciles estimated request cost with observed service. Separate counters
avoid inventing a conversion between CPU preparation seconds and GPU inference
service. This design extends the fairness boundary from model-ready admission
back to request arrival without modifying the model, video decoder, or GPU
scheduler.
\end{abstract}

\maketitle

\section{Introduction}

Large multimodal models extend language-model serving to images, audio, and
video, enabling applications such as video question answering, content search,
assistive analysis, and long-form media understanding. This capability adds a
substantial execution path before and alongside language-model inference.
Serving a video request may require fetching the media, parsing its container,
seeking to temporal positions, decoding and resizing frames, constructing
visual inputs, running a vision encoder, prefilling the language model, and
generating output tokens. Recent studies show that these additional stages can
materially affect time to first token, memory use, throughput, and resource
utilization~\cite{qiu2025efficient,singh2025epd,zhao2025flashcodec}.

Production services multiplex this pipeline among multiple clients or
tenants. A useful shared service must provide more than high aggregate
throughput: it must isolate clients from one another. Without isolation, a
tenant submitting many requests, or a small number of unusually expensive
videos, can occupy preparation workers and delay every other tenant. This
problem is not addressed by equal request counts. One request may sample eight
frames from an easily seekable video, while another may decode 128 frames from
a long, high-resolution stream with sparse keyframes. Giving both tenants the
same number of completed requests can give them very different shares of the
underlying resources.

Fairness is important for three related reasons. First, it contains aggressive
or misconfigured clients, so load from one tenant does not produce unbounded
interference for others. Second, it provides a meaningful service contract:
continuously backlogged tenants receive comparable shares of contended service
rather than comparable request counts. Third, it preserves work conservation.
When some tenants are idle, active tenants should use the spare capacity rather
than being constrained by static per-client reservations or rate limits.

Fair scheduling has recently been studied for text-only LLM inference. Virtual
Token Counter (VTC) defines service using a cost function over processed input
and output tokens and selects clients according to their accumulated virtual
service~\cite{sheng2024vtc}. It integrates fairness with continuous batching,
where requests can join a running batch as other requests finish. This is a
significant improvement over requests-per-minute limits because it reflects
heterogeneous inference cost and remains work-conserving. Systems such as Orca
and vLLM similarly demonstrate why high-throughput generation relies on
iteration-level or continuous batching rather than serial request
execution~\cite{yu2022orca,kwon2023vllm}.

Inference fairness, however, begins too late for multimodal requests. A video
request is invisible to an inference scheduler until its media has been
prepared and its model input is ready. Before that point, it may wait for and
consume CPU cores, decoder contexts, storage bandwidth, memory bandwidth, or a
GPU video decoder. A tenant can therefore look underserved to an inference
scheduler while it is simultaneously consuming a disproportionate fraction of
upstream preparation capacity. Conversely, two tenants can receive similar
token service inside the model engine after having experienced substantially
different preparation delays.

This boundary mismatch means that inference fairness does not compose into
end-to-end multimodal fairness. Consider two continuously backlogged tenants.
Tenant $A$ submits inexpensive videos, while tenant $B$ submits videos with
larger frame budgets and more expensive seeking and decoding. An
inference-only scheduler can equalize their model service only after requests
become model-ready. It cannot prevent $B$ from occupying the preparation tier,
nor can it compensate $A$ for time spent waiting upstream. Adding faster GPUs
does not resolve contention at the preparation tier, and adding preparation
workers does not define how that capacity should be allocated under overload.

Recent multimodal systems expose precisely why an end-to-end treatment is
necessary. EPD disaggregation separates encoding, prefill, and decoding to
improve resource allocation and interference management~\cite{singh2025epd}.
FlashCodec accelerates video preprocessing through collaborative GPU decoding,
while UnifiedServe coordinates vision and language-model stages through
logical decoupling and physical resource sharing~\cite{zhao2025flashcodec}.
These systems optimize stage placement, decoding latency, interference, and
aggregate throughput. Their request-level decisions do not define how multiple
tenants should share both media preparation and model inference. Conversely,
VTC accounts for client service inside LLM inference but not the media work
consumed before engine admission. The two lines of work leave an important
gap: selecting the next waiting multimodal request according to tenant service
across the complete serving path.

We address this gap with stage-wise, unweighted max-min fairness. Whenever
capacity becomes available at either preparation or inference admission, our
scheduler selects the backlogged tenant with the least accumulated service at
that stage and dispatches that tenant's oldest request. A costly video is not
rejected or assigned a lower request class. Instead, its tenant is charged for
the service it consumes and is selected less often until other continuously
backlogged tenants catch up. Fairness therefore means equal resource service,
not equal request completions or equal latency for heterogeneous requests.

The design maintains separate preparation and inference counters. This choice
is deliberate: CPU decoding time and GPU model service have no principled
fixed exchange rate. Preparation cost depends on codec, keyframe placement,
resolution, duration, frame selection, and host contention, whereas inference
cost depends on visual tokens, prompt and output length, and batch composition.
We estimate each request's stage cost at dispatch to avoid admitting many
expensive requests concurrently, then reconcile the estimate with observed
service after completion. The scheduler remains external to the model, media
decoder, and GPU scheduler.

This paper makes the following contributions:

\begin{itemize}
    \item It identifies the fairness boundary mismatch between multimodal
    preparation and inference-only fair scheduling.
    \item It formulates multimodal serving as unweighted max-min allocation at
    two asynchronous stages with different service units.
    \item It designs and implements a work-conserving scheduler with online
    cost estimation and observed-cost reconciliation at both stages.
\end{itemize}

\section{Multimodal Serving Model}

Let a request be
\begin{equation}
r = \langle id_r, u_r, a_r, v_r, x_r, f_r \rangle,
\end{equation}
where $u_r$ is the authenticated tenant, $a_r$ is the arrival time, $v_r$ is
the video, $x_r$ is the text query, and $f_r$ is the requested frame budget.
The fairness policy does not require an application-assigned request class.
All requests belonging to a tenant share the same entitlement, and requests
within the selected tenant are ordered by arrival time.

A request passes through two contended stages:

\begin{enumerate}
    \item \textbf{Preparation ($P$):} media access, metadata parsing, frame
    selection, seeking, decoding, resizing, and prompt construction.
    \item \textbf{Inference ($I$):} vision encoding, language-model prefill,
    and autoregressive generation after model-ready admission.
\end{enumerate}

Preparation and inference overlap across requests, but each has a finite
admission capacity and its own waiting queue. End-to-end latency is
\begin{equation}
T_{\mathrm{E2E}}(r) =
T^P_q(r) + T^P_s(r) + T^I_q(r) + T^I_s(r),
\end{equation}
where $T^s_q$ and $T^s_s$ are queueing and service time at stage $s$.
Inference-only fairness controls only the final two terms.

\section{Existing Fairness Approaches}

\paragraph{Per-client rate limits.}
Requests-per-minute, token-rate, and concurrency limits bound admission over a
time window. They protect against overload but treat heterogeneous requests as
equal and may leave capacity idle when some tenants are inactive. They do not
determine which admitted request should receive the next preparation slot.

\paragraph{Classical fair queueing.}
Weighted Fair Queueing, Start-Time Fair Queueing, and Deficit Round Robin
approximate fluid sharing by scheduling flows according to virtual service.
They generally assume a known packet size, a single shared resource, or a
preemptible service quantum. Multimodal requests instead consume uncertain
amounts of service at multiple asynchronous stages.

\paragraph{Max-min fairness.}
Max-min fairness maximizes the smallest allocation among backlogged tenants;
after that minimum is fixed, it maximizes the next smallest allocation, and so
on. With equal tenant weights and equal demands, each backlogged tenant
receives an equal share. If a tenant cannot use its share, the unused capacity
is redistributed among active tenants, making the policy work-conserving.

\paragraph{VTC and inference fairness.}
VTC adapts fair queueing to LLM inference by maintaining a virtual token
counter per client and charging processed input and output tokens using a
customizable cost function~\cite{sheng2024vtc}. It addresses variable prompt and generation
lengths within the inference engine. It does not charge video fetching,
decoding, or transformation performed before engine admission.

\paragraph{Multi-resource fairness.}
Dominant Resource Fairness reasons about simultaneous shares of resources such
as CPU and memory. Our pipeline differs because preparation and inference are
ordered, asynchronously pipelined stages. A request cannot consume inference
service until preparation completes. We therefore avoid treating CPU and GPU
service as a simultaneous resource vector or combining them into one unit.

\section{Why Existing Methods Do Not Apply Directly}

\paragraph{There is no common service unit.}
Preparation cost depends on codec, keyframe placement, duration, resolution,
frame budget, storage behavior, and host contention. Inference cost depends on
visual-token count, prompt length, output length, and batch composition. One
CPU second is not naturally equivalent to a fixed number of model tokens.

\paragraph{Request cost is initially uncertain.}
The exact preparation cost is unknown before decoding, and output length is
unknown before generation finishes. Algorithms that require an exact finish
tag before dispatch cannot be used unchanged.

\paragraph{Requests become visible at different times.}
The inference scheduler cannot select a compressed video that is still being
prepared. Equal GPU service therefore says nothing about how CPU preparation
capacity was allocated.

\paragraph{Preparation is coarse-grained and non-preemptive.}
Interrupting active decoding may discard useful work and require decoder
restart. Fairness decisions occur when a worker becomes available, so a new
request can overtake queued work but cannot immediately displace active work.

\paragraph{Continuous batching changes inference capacity.}
The effective cost and throughput of inference depend on the current batch.
Long prompts, visual tokens, and sequential output decoding do not consume
capacity in the same way. A fixed requests-per-second entitlement is therefore
not a stable measure of inference service.

\section{Fairness Objective}

Let $\mathcal{B}_s(t)$ be the set of tenants backlogged at stage
$s\in\{P,I\}$ at time $t$. Let $W^s_u(t_1,t_2)$ be the observed service
received by tenant $u$ at stage $s$ during $[t_1,t_2]$. We seek unweighted
max-min fairness independently at each stage. In an ideal fluid server, two
continuously backlogged tenants receive equal service:
\begin{equation}
W^s_u(t_1,t_2) = W^s_v(t_1,t_2),
\quad u,v\in\mathcal{B}_s(t),\;t\in[t_1,t_2].
\end{equation}

With indivisible, non-preemptive requests, exact equality at every instant is
impossible. The practical objective is bounded discrepancy:
\begin{equation}
\left|W^s_u(t_1,t_2)-W^s_v(t_1,t_2)\right| \leq \Delta_s,
\label{eq:bounded-fairness}
\end{equation}
where $\Delta_s$ depends on stage concurrency, the largest non-preemptive
service unit, and cost-estimation error, but should not grow with
$t_2-t_1$.

The scheduler must additionally be work-conserving:
\begin{equation}
\mathcal{B}_s(t)\neq\emptyset
\Longrightarrow
\text{no admissible slot at stage $s$ remains idle}.
\end{equation}

This definition equalizes consumed resource service, not completed request
counts and not per-request latency. If tenant $A$ submits more expensive videos
than tenant $B$, $A$ may complete fewer requests while receiving the same
preparation share. This is intentional: allowing both tenants to complete the
same number of requests would give $A$ more resource service.

\section{Design}
\label{sec:design}

Our scheduler is an external frontend to the media-preparation workers and
inference engine. It does not modify the model, video decoder, CUDA scheduler,
or continuous-batching implementation. Instead, it controls which waiting
request is admitted when capacity becomes available at either stage.

\subsection{Separate Stage Counters}

For every tenant $u$, the scheduler maintains
\begin{equation}
V^P_u \quad\text{and}\quad V^I_u,
\end{equation}
the accumulated virtual service at preparation and inference. Maintaining
separate counters is essential: summing CPU seconds and GPU tokens would embed
an arbitrary exchange rate and could hide unfairness at one stage behind
service at the other.

At stage $s$, the next tenant is
\begin{equation}
u^* = \arg\min_{u:Q^s_u\neq\emptyset} V^s_u.
\label{eq:tenant-selection}
\end{equation}
The scheduler then selects the oldest request in $Q^s_{u^*}$. Tenant identity
is retained across both stages, but each stage makes its selection from its
own backlogged set and counter.

\subsection{In-Flight Service Accounting}

Max-min selection depends on an up-to-date view of the service received by
each tenant. However, the true cost of a request becomes available only after
the request completes. This delay matters when preparation workers or
inference slots operate concurrently. If a running request contributes no
service until completion, its tenant can continue to appear least served and
may be assigned several additional slots before any of its work is reflected
in the counter.

Waiting until completion to charge service can dispatch several expensive
requests from the same tenant concurrently. We therefore charge predicted
cost $\widehat{c}_s(r)$ at dispatch:
\begin{equation}
V^s_{u(r)} \leftarrow V^s_{u(r)} + \widehat{c}_s(r).
\end{equation}
When the request completes, the prediction is replaced with observed service:
\begin{equation}
V^s_{u(r)} \leftarrow V^s_{u(r)}
  + c_s(r)-\widehat{c}_s(r).
\label{eq:reconcile}
\end{equation}

The preparation estimator groups observations using available metadata such
as backend, frame count, duration bucket, resolution, and codec, and updates
an exponentially weighted moving average. The inference estimator learns from
observed request service and may use frame or visual-token count as a feature.
Prediction affects short-term dispatch order, while reconciliation ensures
that persistent accounting reflects observed rather than predicted cost.
Thus, the estimate acts only as a reservation for unfinished work; fairness
accounting ultimately uses measured service. With a single serial worker this
reservation is unnecessary, but it prevents transient monopolization when
multiple requests can be in flight.

\begin{algorithm}[t]
\caption{Stage-wise Unweighted Max-Min Scheduling}
\label{alg:max-min}
\begin{algorithmic}[1]
\Require stage $s$, tenant queues $Q^s_u$, counters $V^s_u$
\When{a slot becomes available at stage $s$}
    \State $A \gets \{u\mid Q^s_u\neq\emptyset\}$
    \If{$A\neq\emptyset$}
        \State $u^* \gets \arg\min_{u\in A}(V^s_u,u)$
        \State $r^* \gets \arg\min_{r\in Q^s_{u^*}}(a_r,id_r)$
        \State $\widehat c \gets \Call{PredictCost}{s,r^*}$
        \State $V^s_{u^*}\gets V^s_{u^*}+\widehat c$
        \State dispatch $r^*$ at stage $s$
    \EndIf
\EndWhen
\When{$r$ completes at stage $s$ with observed cost $c$}
    \State $V^s_{u(r)}\gets V^s_{u(r)}+c-\widehat c_s(r)$
    \State \Call{UpdateCostModel}{$s,r,c$}
\EndWhen
\end{algorithmic}
\end{algorithm}

\subsection{Tenant Activation}

A newly active tenant should join at the current fair-service frontier rather
than at a historical zero. Otherwise it could monopolize the server until its
counter catches up with long-running tenants. On activation, its stage counter
is initialized to at least the minimum counter among active tenants. An idle
tenant does not accumulate credit, but unused capacity is immediately consumed
by tenants that remain backlogged.

\subsection{Bounded Admission}

Preparation and inference use bounded concurrency. A bounded model-ready queue
prevents one stage from producing an arbitrarily large backlog that fixes the
downstream order too early. At every newly available slot, the scheduler can
reconsider the least-served tenant using current counters and backlog state.

\section{Evaluation Methodology}

The main evaluation compares three policies on identical arrival traces:

\begin{itemize}
    \item \textbf{FCFS:} neither stage performs tenant-fair selection;
    \item \textbf{Inference-only fairness:} preparation is FCFS, while
    model-ready requests are ordered using per-tenant inference service; and
    \item \textbf{Cross-stage max-min:} both preparation and inference select
    the least-served backlogged tenant using separate service counters.
\end{itemize}

Traces should contain at least three tenants and heterogeneous media costs,
including different frame budgets, codecs, durations, and resolutions. All
policies must receive the same requests, videos, arrival times, random seeds,
worker counts, and inference concurrency.

We report the following metrics separately for preparation and inference:

\begin{itemize}
    \item service received by each tenant during continuously backlogged
    intervals;
    \item maximum pairwise service discrepancy and whether it grows with the
    interval length;
    \item Jain's fairness index over tenant service and tenant slowdown;
    \item per-tenant queueing time, end-to-end latency, and throughput;
    \item aggregate throughput and resource utilization to establish work
    conservation; and
    \item prediction error before and after online cost-model adaptation.
\end{itemize}

The key comparison is inference-only fairness versus cross-stage max-min.
This isolates whether extending accounting to media preparation improves
end-to-end isolation without materially reducing aggregate throughput.

\section{Expected Claims and Non-Claims}

The intended claim is that stage-wise service accounting prevents a tenant
from obtaining an unbounded advantage at either preparation or inference while
the tenants remain backlogged. The system aims to preserve work conservation
and reduce cross-tenant disparity in end-to-end slowdown.

The design does not claim that all requests have equal latency, that all
tenants complete equal request counts, or that CPU and GPU service are
interchangeable. It also does not provide instantaneous fairness when a large
non-preemptive request is already executing. These distinctions are important
for interpreting both the formulation and the measurements.

\section{Conclusion}

Fair LLM inference is not sufficient for fair multimodal serving. Video
requests consume heterogeneous and contended service before the inference
engine can observe them. We extend the fairness boundary to request arrival by
applying unweighted max-min selection independently at preparation and
inference admission. Separate virtual-service counters, online estimation,
observed-cost reconciliation, and work-conserving dispatch make the approach
compatible with an asynchronous CPU--GPU pipeline without modifying its media
decoder or inference engine.

\bibliographystyle{ACM-Reference-Format}
\bibliography{multimodal_max_min_fairness}

\end{document}
